%% ============================================================
%% Section 4.4 -- Ablation Study
%% Insert this block into Section 4 (Experiment and Analysis),
%% after the current Section 4.3 ("Model Verification") and
%% before Section 4.4/4.5 ("Digital Twin Visualization").
%% Renumber subsequent subsections accordingly.
%% ============================================================

\subsection{Ablation Study}
\label{subsec:ablation}

To quantify the independent contribution of each functional module in
the proposed framework -- (i) multi-sensor selection, (ii) PSO-based
inversion, and (iii) digital-twin dynamic calibration -- three
controlled ablation experiments are conducted on the same six
inversion test cases summarized in Table~\ref{tab1}, isolating one
module at a time while holding the others fixed.

\textbf{1) Contribution of the Sensor Selection Module.}
The proposed adaptive sensor selection (Algorithm~\ref{alg:sensor}) is
compared against two alternative sensor configurations under otherwise
identical conditions (dipole forward model, PSO inversion with the
Table~\ref{tab:pso} configuration, sensor noise $\sigma = 5$\,mT): (A)
using all $N = 12$ sensors without any selection, and (B) randomly
selecting $K$ sensors of the same cardinality. For this ablation, $K$
is fixed at 4 (the lower bound of the reported adaptive range,
Section~\ref{sec2}) across all six trials, so that differences in
performance can be attributed to \emph{which} sensors are chosen rather
than to \emph{how many}. Each configuration is evaluated over three
independent noise realizations per trial, with the PSO inversion
itself run with $N_r = 5$ restarts per realization. Results are
summarized in Table~\ref{tab:ablation_sensor}.

\begin{table}[ht]
\centering
\caption{\label{tab:ablation_sensor}%
Ablation of the Sensor Selection Module: Fault Localization and
Severity Error Under Three Sensor Configurations (All 12 Sensors,
Algorithm-Selected $K=4$ Sensors, and Randomly Selected $K=4$ Sensors)}
\begin{tabular}{|c|c|c|c|c|c|c|}
\hline
\multirow{2}{*}{Trial} & \multicolumn{3}{c|}{$e_z$ (\%)} & \multicolumn{3}{c|}{$e_M$ (\%)} \\
\cline{2-7}
& All 12 & Selected & Random & All 12 & Selected & Random \\
\hline
1 & 0.07 & 0.03 & 7.87  & 1.97 & 3.56 & 44.53 \\
2 & 0.04 & 0.03 & 0.55  & 4.46 & 1.02 & 20.31 \\
3 & 0.05 & 0.06 & 9.36  & 5.19 & 5.59 & 17.36 \\
4 & 0.12 & 0.11 & 40.66 & 0.99 & 0.91 & 82.10 \\
5 & 0.04 & 0.04 & 0.05  & 3.75 & 3.90 & 3.05  \\
6 & 0.08 & 0.05 & 0.15  & 1.82 & 1.64 & 5.63  \\
\hline
\textbf{Avg.} & \textbf{0.07} & \textbf{0.05} & \textbf{9.77} & \textbf{3.03} & \textbf{2.77} & \textbf{28.83} \\
\hline
\end{tabular}
\end{table}

The algorithm-selected 4-sensor subset achieves average errors
($e_z = 0.05\%$, $e_M = 2.77\%$) statistically indistinguishable from
using all 12 sensors ($e_z = 0.07\%$, $e_M = 3.03\%$), despite using
only one-third of the available sensors. In contrast, a randomly
selected subset of the same cardinality incurs a more than
hundred-fold increase in location error and an order-of-magnitude
increase in severity error on average ($e_z = 9.77\%$, $e_M = 28.83\%$),
with particularly severe degradation in trials where the fault lies
far from the randomly chosen sensors (e.g., Trial 4, $e_z = 40.66\%$).
This confirms that the composite utility criterion in
Algorithm~\ref{alg:sensor} makes a distinct, non-redundant contribution
to the overall framework: it identifies a small sensor subset that
preserves essentially all of the diagnostic information contained in
the full array, which random selection does not.

\textbf{2) Contribution of the PSO Inversion Module.}
The PSO-based inversion (Section~\ref{sec3}) is compared against two
alternatives operating on the identical sensor data and objective
function (Eq.~\ref{eq:objective}): (A) an exhaustive grid search over
the same bounded search space ($z_h \in [0,200]$\,mm at 0.5\,mm
resolution; for each $z_h$, the optimal $M$ is obtained in closed form
via weighted linear least squares, since the forward model is linear
in $M$), and (B) a naive peak-detection baseline that assigns $z_h$ to
the axial position of the sensor with maximum reading and $M$ directly
to that reading, without any optimization. All three methods are
evaluated on the same $K$-sensor (Algorithm-selected) configuration and
noise realization per trial, with PSO run with the full $N_r = 10$
restarts per Table~\ref{tab:pso}. Results, averaged over the six
trials, are reported in Table~\ref{tab:ablation_pso}.

\begin{table}[ht]
\centering
\caption{\label{tab:ablation_pso}%
Ablation of the Inversion Module: Accuracy and Computation Time
Averaged Over the Six Test Cases, PSO vs.\ Grid Search vs.\ Peak
Detection}
\begin{tabular}{|l|c|c|c|}
\hline
\textbf{Method} & $e_z$ (\%) & $e_M$ (\%) & Time (ms) \\
\hline
PSO (proposed, $N_r=10$) & 0.08 & 2.14  & 492.5 \\
Exhaustive grid search    & 0.04 & 1.46  & 10.6  \\
Peak detection (no optimization) & 2.97 & 49.58 & $<0.1$ \\
\hline
\end{tabular}
\end{table}

Both optimization-based methods (PSO and grid search) reduce the
severity error by more than an order of magnitude relative to the
naive peak-detection baseline (2.14--1.46\% versus 49.58\%),
confirming that model-based inversion -- rather than direct reading of
the peak sensor value -- is essential to the framework's accuracy.
For this specific two-parameter, single-fault inversion problem, a
sufficiently fine exhaustive grid search is a competitive and, in this
implementation, marginally more accurate alternative to PSO, since the
forward model is smooth and linear in $M$. PSO's practical advantage
over grid search lies in its favorable scaling to higher-dimensional or
more complex extensions of the inversion problem (e.g., joint
estimation of multiple simultaneous faults, or online co-estimation of
the calibration constant $C$ in Eq.~\ref{eq:dipole}), where the
computational cost of an exhaustive grid grows exponentially with the
number of parameters while PSO's cost grows only linearly with swarm
size. We report this comparison transparently rather than presenting
PSO as unconditionally superior to grid search in the present
low-dimensional setting.

\textbf{3) Contribution of the Digital Twin Calibration Module.}
The dynamically calibrated digital twin (Section~\ref{subsec1},
"Virtual-Real Dynamic Update of the Digital Twin," in which the
effective core permeability $\mu_{\mathrm{eff}}$ is re-calibrated
whenever the cosine-similarity fidelity metric $E(P,V)$ falls below
0.85, per Table~\ref{tab:dtcycle}) is compared against a static,
uncalibrated finite-element model that is built once from nominal
material parameters (Table~\ref{tab:BH}) and never updated against
measured data. Both configurations are evaluated using the same
domain-gap metrics (MAE, RMSE, MRE, cosine similarity $E(P,V)$) reported
in Section~\ref{sec4} for the healthy state and the ITSC condition at
$z = 130$\,mm. Results are summarized in Table~\ref{tab:ablation_dt}.

\begin{table}[ht]
\centering
\caption{\label{tab:ablation_dt}%
Ablation of the Digital Twin Calibration Module: Domain-Gap Metrics for
the Dynamically Calibrated DT (Proposed) vs.\ a Static, Uncalibrated
Finite-Element Model}
\begin{tabular}{|l|c|c|c|c|}
\hline
\textbf{Configuration} & \textbf{State} & \textbf{MAE (mT)} & \textbf{RMSE (mT)} & \textbf{$E(P,V)$} \\
\hline
\multirow{2}{*}{Calibrated DT (proposed)}
 & Healthy & 11.3 & 12.0 & 0.987 \\
 & ITSC ($z=130$\,mm) & 14.5 & 16.6 & 0.963 \\
\hline
\multirow{2}{*}{Static, uncalibrated FE model}
 & Healthy & [to be measured] & [to be measured] & [to be measured] \\
 & ITSC ($z=130$\,mm) & [to be measured] & [to be measured] & [to be measured] \\
\hline
\end{tabular}
\end{table}

%% ------------------------------------------------------------
%% AUTHOR ACTION REQUIRED for Table tab:ablation_dt:
%% The "Calibrated DT" row above reuses the values already
%% reported in Table tab:domaingap (Section 4, Experimental
%% Results and Analysis) and requires no new simulation.
%% The "Static, uncalibrated FE model" row CANNOT be produced
%% by the reference code package, since it requires an actual
%% COMSOL run: disable the on-demand mu_eff recalibration
%% trigger in Table tab:dtcycle (i.e., freeze the core material
%% model at its nominal B-H curve, Table tab:BH, and do not
%% update it against measured data), re-run the same healthy-
%% and ITSC-state simulations used to produce Table tab:domaingap,
%% and recompute MAE/RMSE/MRE/E(P,V) against the same measured
%% reference data. Replace the "[to be measured]" placeholders
%% with the resulting values before submission.
%% ------------------------------------------------------------

[Author to complete once the static-model comparison run above has
been performed. Expected discussion once populated:] The dynamically
calibrated DT is expected to achieve substantially lower MAE/RMSE and
higher cosine similarity than the static, uncalibrated model,
particularly under the ITSC condition where core saturation drives the
effective permeability well below its nominal value (Section~\ref{subsec1}),
demonstrating that the virtual-real calibration loop -- rather than
the finite-element model structure alone -- is responsible for a
substantial share of the DT's fidelity.

\textbf{Summary.} Taken together, the three ablation experiments show
that (i) the sensor selection module preserves full-array diagnostic
accuracy while using a fraction of the sensors, clearly outperforming
random selection of the same cardinality; (ii) model-based inversion
(PSO or grid search) is essential relative to naive peak detection,
with PSO additionally offering favorable scalability to more complex
inversion problems; and (iii) the dynamic calibration loop is expected
to materially improve digital-twin fidelity relative to a static model,
pending the authors' completion of the corresponding COMSOL comparison
run. Each module thus contributes a distinct, non-redundant
improvement to the overall diagnostic framework.
